% Nguồn SGK Toán 9 Tập hai — Bài tập cuối chương VIII, trang 66:
% https://cdn3.olm.vn/upload/thuvienso/4491669493-page-66-1771844556377.png
%
% Nguồn SBT Toán 9 Tập hai — Ôn tập chương VIII, PDF page 47--48:
% SBT TOÁN 9 TẬP 2.pdf
% SHA-256: 5ff01fd213d1adc75f9c54b4408d6a9642a4d8038e685040efef3311a196196c
% PDF page 47 là trang chồng ranh giới: chỉ lấy phần bắt đầu từ ÔN TẬP CHƯƠNG VIII.
% Không lấy phần lời giải/hướng dẫn/đáp số của SBT.

\section*{Bài tập cuối chương VIII}

\begin{exercises}
    \subsection{Bài tập cuối chương VIII}

    \textbf{A. TRẮC NGHIỆM}

    \begin{questions}[resume=SGK]
        \item Gieo đồng thời hai con xúc xắc cân đối, đồng chất. Xác suất để “Tổng số chấm xuất hiện trên hai con xúc xắc lớn hơn hoặc bằng 10” là
        \begin{tasks}(4)
            \task[(A)] $\dfrac{7}{36}$.
            \task[(B)] $\dfrac{2}{9}$.
            \task[(C)] $\dfrac{1}{6}$.
            \task[(D)] $\dfrac{5}{36}$.
        \end{tasks}
        \par\noindent\answerbox\par

        \item Có hai túi I và II. Túi I chứa 4 tấm thẻ, đánh số 1; 2; 3; 4. Túi II chứa 5 tấm thẻ, đánh số 1; 2; 3; 4; 5. Rút ngẫu nhiên một tấm thẻ từ mỗi túi I và II. Xác suất để cả hai tấm thẻ rút ra đều ghi số chẵn là
        \begin{tasks}(4)
            \task[(A)] $\dfrac{1}{5}$.
            \task[(B)] $\dfrac{3}{20}$.
            \task[(C)] $\dfrac{1}{4}$.
            \task[(D)] $\dfrac{4}{21}$.
        \end{tasks}
        \par\noindent\answerbox\par

        \item Một túi đựng 4 viên bi có cùng khối lượng và kích thước, được đánh số 1; 2; 3; 4. Lấy ngẫu nhiên hai viên bi từ trong túi. Xác suất để tích hai số ghi trên hai viên bi lớn hơn 3 là
        \begin{tasks}(4)
            \task[(A)] $\dfrac{5}{7}$.
            \task[(B)] $\dfrac{2}{3}$.
            \task[(C)] $\dfrac{3}{4}$.
            \task[(D)] $\dfrac{5}{6}$.
        \end{tasks}
        \par\noindent\answerbox\par
    \end{questions}

    \textbf{B. TỰ LUẬN}

    \begin{questions}[resume=SGK]
        \item Có hai túi I và II. Túi I chứa 3 tấm thẻ, đánh số 2; 3; 4. Túi II chứa 2 tấm thẻ, đánh số 5; 6. Từ mỗi túi I và II, rút ngẫu nhiên một tấm thẻ. Tính xác suất của các biến cố sau:
        \begin{itemize}
            \item $A$: “Hai số ghi trên hai tấm thẻ chênh nhau 2 đơn vị”;
            \item $B$: “Hai số ghi trên hai tấm thẻ chênh nhau lớn hơn 2 đơn vị”;
            \item $C$: “Tích hai số ghi trên hai tấm thẻ là một số chẵn”;
            \item $D$: “Tổng hai số ghi trên hai tấm thẻ là một số nguyên tố”.
        \end{itemize}
        \answerline
        \answerline
        \answerline

        \item Gieo đồng thời hai con xúc xắc cân đối, đồng chất I và II. Tính xác suất của các biến cố sau:
        \begin{itemize}
            \item $E$: “Tổng số chấm xuất hiện trên hai con xúc xắc bằng 11”;
            \item $F$: “Tổng số chấm xuất hiện trên hai con xúc xắc bằng 8 hoặc 9”;
            \item $G$: “Tổng số chấm xuất hiện trên hai con xúc xắc nhỏ hơn 6”.
        \end{itemize}
        \answerline
        \answerline
        \answerline

        \item Hai bạn Minh và Huy chơi một trò chơi như sau: Minh chọn ngẫu nhiên một số trong tập hợp $\{5;\ 6;\ 7;\ 8;\ 9;\ 10\}$; Huy chọn ngẫu nhiên một số trong tập hợp $\{4;\ 5;\ 7;\ 8;\ 9;\ 11\}$. Bạn nào chọn được số lớn hơn sẽ là người thắng cuộc. Nếu hai số chọn được bằng nhau thì kết quả là hoà. Tính xác suất của các biến cố sau:
        \begin{questions}
            \item $A$: “Bạn Minh thắng”.
            \item $B$: “Bạn Huy thắng”.
        \end{questions}
        \answerline
        \answerline
        \answerline
    \end{questions}
\end{exercises}

\section*{Ôn tập chương VIII}

\begin{practice}
    \subsection{Ôn tập chương VIII}

    \textbf{A. Câu hỏi (Trắc nghiệm)}

    \begin{enumerate}
        \item Có hai túi I và II. Túi I chứa ba quả cầu ghi các số 1, 2, 3. Túi II chứa bốn tấm thẻ ghi các số 1, 2, 3, 4. Lấy ngẫu nhiên một quả cầu và một tấm thẻ từ mỗi túi I và II. Xác suất của biến cố “Tích hai số ghi trên quả cầu và tấm thẻ bằng 6” là
        \begin{tasks}(4)
            \task[(A)] $\dfrac{1}{6}$.
            \task[(B)] $\dfrac{1}{7}$.
            \task[(C)] $\dfrac{2}{11}$.
            \task[(D)] $\dfrac{1}{4}$.
        \end{tasks}
        \par\noindent\answerbox\par

        \item Hai bạn Minh và Dũng mỗi người gieo ngẫu nhiên một con xúc xắc cân đối. Xác suất để số chấm xuất hiện trên hai con xúc xắc như nhau là
        \begin{tasks}(4)
            \task[(A)] $\dfrac{2}{11}$.
            \task[(B)] $\dfrac{1}{7}$.
            \task[(C)] $\dfrac{1}{6}$.
            \task[(D)] $\dfrac{1}{4}$.
        \end{tasks}
        \par\noindent\answerbox\par

        \item Bạn Sơn gieo một đồng xu cân đối và bạn Hoà gieo đồng thời hai đồng xu cân đối. Xác suất để có hai đồng xu xuất hiện mặt sấp, một đồng xu xuất hiện mặt ngửa là
        \begin{tasks}(4)
            \task[(A)] $\dfrac{3}{8}$.
            \task[(B)] $\dfrac{3}{10}$.
            \task[(C)] $\dfrac{2}{7}$.
            \task[(D)] $\dfrac{4}{31}$.
        \end{tasks}
        \par\noindent\answerbox\par
    \end{enumerate}

    \textbf{B. BÀI TẬP}

    \begin{questions}[resume=SBT]
        \item Bạn An gieo một con xúc xắc cân đối và bạn Bình gieo một đồng xu cân đối. Tính xác suất của các biến cố sau:
        \begin{itemize}
            \item $E$: “Số chấm xuất hiện trên con xúc xắc là 6 và đồng xu xuất hiện mặt sấp”;
            \item $F$: “Số chấm xuất hiện trên con xúc xắc là số lẻ”;
            \item $G$: “Số chấm xuất hiện trên con xúc xắc là số chẵn và đồng xu xuất hiện mặt sấp”;
            \item $H$: “Số chấm xuất hiện trên con xúc xắc là 5 hoặc đồng xu xuất hiện mặt ngửa”.
        \end{itemize}
        \answerline
        \answerline
        \answerline

        \item Trong một trò chơi, có hai bánh xe; mỗi bánh xe được gắn vào trục quay có mũi tên ở tâm. Bánh xe thứ nhất được chia làm bốn hình quạt như nhau và sơn các màu: trắng, đỏ, xanh, vàng. Bánh xe thứ hai được chia làm ba hình quạt như nhau và sơn các màu: đỏ, xanh, vàng. Người chơi quay hai bánh xe. Người chơi đạt giải nhất nếu hai mũi tên dừng lại ở hai hình quạt màu đỏ, đạt giải nhì nếu hai mũi tên dừng lại ở hai hình quạt cùng màu và đạt giải ba nếu có đúng một mũi tên dừng ở hình quạt màu đỏ. Tính xác suất của các biến cố sau:

        \begin{center}
            \begin{tikzpicture}
                \coordinate (O) at (0,0);
                \draw (O) circle (1.45cm);
                \foreach \a in {45,135,225,315} {
                    \draw (O) -- (\a:1.45cm);
                }
                \node at (90:0.92cm) {Trắng};
                \node at (0:0.92cm) {Đỏ};
                \node at (270:0.92cm) {Xanh};
                \node at (180:0.92cm) {Vàng};
                \draw[-{Stealth[length=2.8mm,width=2mm]}] (O) -- (270:0.95cm);
                \draw (O) circle (1.2pt);
            \end{tikzpicture}
            \qquad
            \begin{tikzpicture}
                \coordinate (O) at (0,0);
                \draw (O) circle (1.45cm);
                \foreach \a in {60,180,300} {
                    \draw (O) -- (\a:1.45cm);
                }
                \node at (120:0.92cm) {Xanh};
                \node at (0:0.92cm) {Đỏ};
                \node at (240:0.92cm) {Vàng};
                \draw[-{Stealth[length=2.8mm,width=2mm]}] (O) -- (270:0.95cm);
                \draw (O) circle (1.2pt);
            \end{tikzpicture}
        \end{center}

        \begin{questions}
            \item $E$: “Người chơi đạt giải nhất”;
            \item $F$: “Người chơi đạt giải nhì”;
            \item $G$: “Người chơi đạt giải ba”.
        \end{questions}
        \answerline
        \answerline
        \answerline

        \item Bạn Tuấn viết ba bức thư cho ba người bạn là An, Bình, Cường và viết tên, địa chỉ của ba người bạn đó lên ba chiếc phong bì. Xếp ngẫu nhiên ba bức thư đó vào ba phong bì.
        \begin{questions}
            \item Mô tả không gian mẫu của phép thử.
            \item Tính xác suất của các biến cố sau:
            \begin{itemize}
                \item $E$: “Có đúng một bức thư đúng địa chỉ”;
                \item $F$: “Cả ba bức thư đúng địa chỉ”;
                \item $G$: “Không có bức thư nào đúng địa chỉ”;
                \item $H$: “Có ít nhất một bức thư đúng địa chỉ”.
            \end{itemize}
        \end{questions}
        \answerline
        \answerline
        \answerline

        \item Một tấm bìa hình tròn được chia làm sáu hình quạt tròn có diện tích bằng nhau, ghi các số 1, 2, 3, 4, 5, 6 và được gắn vào trục quay có mũi tên cố định ở tâm. Quay tấm bìa hai lần. Tính xác suất để mũi tên chỉ vào hai hình quạt tròn không đối xứng nhau qua tâm.

        \begin{center}
            \begin{tikzpicture}
                \coordinate (O) at (0,0);
                \draw (O) circle (1.45cm);
                \foreach \a in {0,60,...,300} {
                    \draw (O) -- (\a:1.45cm);
                }
                \node at (90:0.92cm) {1};
                \node at (30:0.92cm) {2};
                \node at (330:0.92cm) {3};
                \node at (270:0.92cm) {4};
                \node at (210:0.92cm) {5};
                \node at (150:0.92cm) {6};
                \draw[-{Stealth[length=2.8mm,width=2mm]}] (O) -- (270:0.95cm);
                \draw (O) circle (1.2pt);
            \end{tikzpicture}
        \end{center}
        \answerline
        \answerline

        \item Màu hạt của đậu Hà Lan có hai kiểu hình là vàng và xanh. Có hai gene ứng với hai kiểu hình này allele trội A và allele lặn a. Hình dạng hạt của đậu Hà Lan có hai kiểu hình là hạt trơn và hạt nhăn. Có hai gene ứng với hai kiểu hình này allele trội B và allele lặn b. Khi cho lai hai cây đậu Hà Lan, cặp gene của cây con lấy ngẫu nhiên một gene từ cây bố và một gene từ cây mẹ. Giả sử cả cây bố và cây mẹ có kiểu hình “Hạt vàng và trơn”. Cây bố có kiểu gene là (Aa, Bb), cây mẹ có kiểu gene là (Aa, Bb). Tính xác suất để cây con có kiểu hình như cây bố và cây mẹ.
        \answerline
        \answerline
        \answerline
    \end{questions}
\end{practice}
